\section{Archimedean Profiles}
\label{pf:section}

For a pair of complex places, set
\[
 t=(1+a|z|^2)^{-1},\qquad u=(1+a|w|^2)^{-1},\qquad
 g_D(z,w)=t^su^sP(t,u),
\]
where
\[
 s=1.14600585,\qquad a=0.0000128633241758,
\]
and
\[
 P(t,u)=\sum_{i,j=0}^3 B_{ij}b_i^3(t)b_j^3(u),\qquad
 b_i^m(t)=\binom mi t^i(1-t)^{m-i},
\]
with
\begin{equation}\label{pf:bernstein-witness}
 B=\begin{pmatrix}
 1&2.4129438977&2.7721799210&6.6857798061\\
 2.4129438977&3.8532983888&4.6371750119&8.1307168183\\
 2.7721799210&4.6371750119&5.4628959084&9.5512647588\\
 6.6857798061&8.1307168183&9.5512647588&13.8659430610
 \end{pmatrix}.
\end{equation}
All finite decimals are exact rationals. The Bernstein basis is
nonnegative and sums to one, so $P$ is bounded above and below by
positive constants on the real square. We prove the hypotheses of
Proposition~\ref{fw:transfer} for these profiles and derive finite bounds
for their functionals. Section~\ref{sec:certificate} evaluates the
bounds at the fixed rational parameters.

\subsection{The Gaussian and the Pair Convolution}
The compact Gaussian $f_C^p=e^{-a_C|z|^2}$ gives
\begin{equation}\label{pf:gaussian-compact}
 J_C=\log(p/2)+\delta\log(a_C/\pi)-a_C/(2p),\qquad
 a_C=2p\delta
\end{equation}
where the displayed scale maximizes the Gaussian functional. For our
choice of $\delta$, we have $a_C<1$. The slope of the general margin as
a function of the signature
is $-\log\pi-2J_C+J_D$, and depends on the chosen profile.

For the pair calculation, let $\alpha,\alpha'>0$ with
$A=\alpha+\alpha'-1>0$. The Laplace identity
\[
 (1+|z|^2)^{-\alpha}
 =\frac1{\Gamma(\alpha)}\int_0^\infty
 r^{\alpha-1}e^{-r-r|z|^2}\,dr
\]
and Gaussian integration give
\begin{equation}\label{pf:laplace-convolution}
 \int_{\mathbb C}(1+|z|^2)^{-\alpha}
          (1+|z+h|^2)^{-\alpha'}\,dz
 =\frac\pi A\,
 \mathbb E(1+T(1-T)|h|^2)^{-A},
 \qquad T\sim\operatorname{Beta}(\alpha,\alpha').
\end{equation}
After integrating $z$, change the two positive Laplace variables to their
sum $r$ and ratio $T$. Completing the square gives the term
$rT(1-T)|h|^2$, and integration in $r$ gives the Beta expectation.
Tonelli justifies the interchanges, since the integrands are nonnegative.

For $A,B>0$ define
\[
 Q_{AB}(b)=\int_{\mathbb R}
 (1+be^{2v})^{-A}(1+be^{-2v})^{-B}\,dv,\qquad b>0.
\]
The two coordinates in the pair convolution reduce to this integral with
\[
 b=a\sqrt{T(1-T)T'(1-T')}\le a/4.
\]
This follows by translating $v$ until the coefficients of $e^{2v}$ and
$e^{-2v}$ agree, which has Jacobian one. The exponential representation is
\[
 \int_{\mathbb R}e^{-a(Re^{2v}+R'e^{-2v})}\,dv
 =K_0(2a\sqrt{RR'}),\qquad
 K_0(x)=\int_0^\infty e^{-x\cosh t}\,dt,
\]
where the substitution $t=2v+\frac12\log(R/R')$ and evenness cancel
their factors of two.

Splitting the $Q_{AB}$ integral at zero yields
\begin{equation}\label{pf:elementary-hyperbola-error}
 \left|Q_{AB}(b)-\log(1/b)
       +\frac{\psi(A)+\psi(B)+2\gamma_E}{2}\right|
 \le (A+B)b/2.
\end{equation}
Deleting the smaller factor on one half leaves
$\frac12\int_b^\infty(1+x)^{-A}dx/x$. Its difference from its
logarithmic constant lies between zero and $Ab/2$, by
$0\le1-(1+x)^{-A}\le Ax$. On that half the factor removed costs at most
$Bb/2$, since $1-(1+be^{-2v})^{-B}\le Bbe^{-2v}$.
The other half interchanges $A$ and $B$. Thus the total positive error
from the two tails and the total loss from the factors removed are both
at most $(A+B)b/2$, proving the bound. The logarithmic constant follows
from \cite[Equation~(5.9.16)]{DLMF},
\[
 \int_0^1\frac{(1-x)^{A-1}-1}{x}\,dx=-\psi(A)-\gamma_E.
\]
We use $\psi=\Gamma'/\Gamma$ and Euler's constant $\gamma_E$. For later
formulas write
\begin{equation}\label{pf:student-constant}
 C_s=-2\psi(s)+\psi(2s)+(2s-1)^{-1}-\gamma_E.
\end{equation}

\subsection{Endpoint Moments and the Fourier Bound}

The compact Gaussian has finite positive mass and overlap, and every
polynomial moment under its endpoint law is finite. Since $P$ is bounded
above and below by positive constants, the pair mass is finite for
$sp>1$. Formula~\eqref{pf:laplace-convolution} and
\eqref{pf:elementary-hyperbola-error} show that the pair overlap is finite:
the Beta variables have positive parameters and finite logarithmic
moments, and $b\le a/4$. Both integrals are positive. The energy
$-\log g_D=s\log(1+a|z|^2)+s\log(1+a|w|^2)-\log P(t,u)$ is bounded
below and coercive, so its sublevel sets are compact.

Finite endpoint-energy second moments hold for every positive polynomial
Student profile $t^su^sP(t,u)$ when $P$ is bounded above and below by
positive constants on the square and $sp>1$. Choose
$0<\epsilon<\min(s,2s-1)$ and absorb each logarithmic square into a
factor $(1+a|z|^2)^\epsilon$. Formula~\eqref{pf:laplace-convolution}
then has positive Beta parameters and positive hyperbola exponents.
The bound in \eqref{pf:elementary-hyperbola-error} is a constant plus
$|\log b|+b(A+B)/2$, which is integrable against the Beta densities.
This proves the second-moment condition, including $s=1$.

For the Fourier bound, write
$P(t,u)=\sum B_{ij}b_i^m(t)b_j^m(u)$ with $B_{ij}>0$. For a tube width
$0<\rho<1$, put $\omega=(\rho+\rho^2)/(1-\rho^2)$ and
\[
 E_{rs}=\binom mr\binom ms\max|\Delta_1^r\Delta_2^sB|,
 \qquad q_0=\frac{\sum_{r+s>0}E_{rs}\omega^{r+s}}{\min B}.
\]
If $q_0<1$, the Bernstein derivative formula puts $P$ in the right
half-plane throughout the product complex tube. Contour shifting then
gives, with $\beta=sp$,
\[
 |\widehat {g_D^p}(\xi_z,\xi_w)|/A_D
 \le (1-\rho^2)^{-2\beta}(1+q_0)^p
 e^{-2\pi\rho(|\xi_z|+|\xi_w|)/\sqrt a}.
\]
Here the tube is taken after the scaling $x=\sqrt a\,z$ in each
complex coordinate, and its imaginary displacement has Euclidean norm at
most $\rho$. If $D(x,y)=1+\sum_j(x_j+iy_j)^2$, then
\[
 |D(x,y)|\ge (1-\rho^2)(1+|x|^2),\qquad
 \left|D(x,y)^{-1}-(1+|x|^2)^{-1}\right|\le\omega.
\]
The latter follows by bounding the numerator by $2|x|\rho+\rho^2$
and the denominator by $(1-\rho^2)(1+|x|^2)^2$.
The Bernstein derivative identity bounds the coefficient of the
$(r,s)$ Taylor term by $E_{rs}$, uniformly at real $(t,u)\in[0,1]^2$.
Hence the complex value of $P$ differs from its positive real value by
at most $q_0P(t,u)$. We can define $P^p$ holomorphically in the tube and have
$|P_{\rm shifted}|^p\le(1+q_0)^pP(t,u)^p$. The displayed bound on $D$
supplies the other factor $(1-\rho^2)^{-2sp}$. Shift each real
Fourier contour by a vector of length $\rho/\sqrt a$ opposite its
frequency. The shifted majorant is integrable since $sp>1$, so
truncation and dominated convergence justify the shift. This proves the
Fourier envelope. The tube condition $q_0<1$, in addition to positivity
on the real square, gives explicit $M,\sigma$ in
\eqref{geo:poisson-condition}.

With the compact Gaussian \eqref{pf:gaussian-compact}, a sufficient
global choice in \eqref{geo:poisson-condition} is
\begin{equation}\label{pf:polynomial-fourier-contract}
 M=\max\{2,\sqrt{K_D}\},\quad
 \sigma=\min\{1,2\pi\rho/\sqrt a\},\quad
 K_D=(1-\rho^2)^{-2sp}(1+q_0)^p.
\end{equation}
There are $b$ one-coordinate blocks and $c$ two-coordinate blocks, so their
prefactors are bounded by $M^{b+2c}=M^d$. The maximum and minimum give
a bound uniform in the signature.

\subsection{Finite Bounds for Polynomial Profiles}

Let $s\ge1$, $\beta=sp>1$, $\eta_0=2s-1$, and
\[
 g_D(z,w)=t^s u^s P(t,u),\qquad
 P(t,u)=\sum_{i,j}d_{ij}t^iu^j>0\quad(0\le t,u\le1).
\]
Put
\[
 A_{ik}=\eta_0+i+k,\quad B_{jl}=\eta_0+j+l,\quad
 w_{ijkl}=\frac{d_{ij}d_{kl}}{A_{ik}B_{jl}},\quad b_0=\sum w_{ijkl}>0,
\]
\[
 c(x,y)=\frac1{x+y-1}
 +\frac{\psi(x+y-1)-\psi(x)-\psi(y)-\gamma_E}{2},
\]
\[
 \kappa_P=\frac{\sum w_{ijkl}[c(s+i,s+k)+c(s+j,s+l)]}{b_0}.
\]
Here $\pi^2b_0=\int_{\mathbb C^2}g_D^2$ at scale $a=1$.
Let $z_0=a^2/16$ and assume, for every cross with nonzero coefficient,
\begin{equation}\label{pf:cross-threshold}
 A_{ik}B_{jl}z_0<1,\qquad
 2\log(4/a)\ge H_{\lceil A_{ik}\rceil-1}+H_{\lceil B_{jl}\rceil-1},
 \qquad \log(4/a)>1/2,
\end{equation}
where $H_0=0$. Define
\[
 E_- =\log(4/a)\sum_{w_{ijkl}<0}(-w_{ijkl})
                \frac{A_{ik}B_{jl}z_0}{1-A_{ik}B_{jl}z_0}.
\]
Then
\begin{equation}\label{pf:polynomial-overlap}
 I_D\ge\pi^2a^{-2}b_0
       [\log(1/a)+\kappa_P-E_-/b_0].
\end{equation}
The bracket must be positive before taking its logarithm.

To prove this bound, recall that for $A,B\ge1$,
\[
 Q_{AB}(b)=\int_{\mathbb R}(1+be^{2u})^{-A}(1+be^{-2u})^{-B}\,du.
\]
For $0<b<1$, its convergent logarithmic series is
\[
 Q_{AB}(b)=\sum_{n\ge0}\frac{(A)_n(B)_n}{(n!)^2}b^{2n}
 \left[\log(1/b)+\psi(n+1)-\frac{\psi(A+n)+\psi(B+n)}2\right].
\]
To verify the series, substitute $x=be^{2u}$ and then $t=x/(x+b^2)$.
This expresses $2Q_{AB}(\sqrt z)$ as
\[
 \int_0^1t^{B-1}(1-t)^{A-1}(1-(1-z)t)^{-A}\,dt.
\]
Differentiation and integration by parts give
$z(1-z)y''+[1-(A+B+1)z]y'-AB y=0$.
Substituting the displayed series verifies its coefficients, using
$a_{n+1}/a_n=(n+A)(n+B)/(n+1)^2$ and the digamma recurrence.
The coefficients have polynomial growth, so the series and its fixed
derivatives converge locally uniformly for $0<z<1$. Its leading
logarithm and constant agree with the elementary estimate~\eqref{pf:elementary-hyperbola-error}. To see uniqueness with these two leading terms, let
$\phi(z)=\sum a_nz^n$, so $\phi(0)=1$ and $\phi>0$ on $[0,1)$.
For the difference $v$ of two such solutions, reduction of order gives
$(v/\phi)'=C/[z(1-z)^{A+B}\phi^2]$. Thus
$v/\phi=C\log z+D+o(1)$, and the matching leading terms force $C=D=0$.

The digamma recurrence and monotonicity imply
\[
 -H_{\lceil A\rceil-1}-H_{\lceil B\rceil-1}
 \le2\psi(n+1)-\psi(A+n)-\psi(B+n)\le0.
\]
Under the threshold in \eqref{pf:cross-threshold}, every term with
$n\ge1$ is nonnegative. Since $(A)_n(B)_n/(n!)^2\le(AB)^n$, summation gives
\[
 0\le Q_{AB}(b)-\log(1/b)
           +\tfrac12[\psi(A)+\psi(B)+2\gamma_E]
 \le\frac{ABb^2\log(1/b)}{1-ABb^2}.
\]
The right side increases for $b\le a/4$ under the stated conditions.
Apply this inside the two Beta expectations from \eqref{pf:laplace-convolution}, where $b=a\sqrt{T(1-T)T'(1-T')}\le a/4$.
Their logarithmic expectations give $c(x,y)$. For positive cross
coefficients we discard the nonnegative remainder. For negative cross
coefficients we subtract its upper bound. This proves
\eqref{pf:polynomial-overlap}, allowing either sign of monomial
coefficient.

To retain cancellation between monomial coefficients, we collect exact
terms before bounding the signed remainder. For $n\ge0$, define
\[
 K_n(i,k)=\frac{(s+i)_n(s+k)_n}
 {n!(\eta_0+i+k+n)_{n+1}},\qquad
 S_x(j)=\sum_{h=0}^{j-1}\frac1{x+h},
\]
with $S_x(0)=0$, and put
\[
 \begin{split}
 r_n(i,k)={}&-\eta_0^{-1}+H_n/2
 -S_{\eta_0}(i+k+n)/2+S_{\eta_0}(i+k+2n+1)\\
 &-[S_s(i+n)+S_s(k+n)]/2,\qquad R_n(i,k)=K_n(i,k)r_n(i,k),
 \end{split}
\]
\[
 B_n=\sum d_{ij}d_{kl}K_n(i,k)K_n(j,l),\qquad
 V_n=\sum d_{ij}d_{kl}
 [R_n(i,k)K_n(j,l)+K_n(i,k)R_n(j,l)].
\]
For rational parameters these are rational numbers. Under
\eqref{pf:cross-threshold}, for every integer $m\ge0$,
\begin{equation}\label{pf:collected-overlap}
 \begin{split}
 \frac{I_D}{\pi^2a^{-2}}\ge{}&
 \sum_{n=0}^m a^{2n}\{B_n[\log(1/a)+C_s]+V_n\}-E_{-,m},\\
 E_{-,m}={}&\log(4/a)\sum_{w_{ijkl}<0}(-w_{ijkl})
 \frac{(A_{ik}B_{jl}z_0)^{m+1}}{1-A_{ik}B_{jl}z_0}.
 \end{split}
\end{equation}
Thus $m=1$ retains the exact $a^2$ term and leaves a signed error of order
$a^4\log(1/a)$ for a fixed polynomial. The displayed right-hand side is a
lower bound for the overlap and gives no upper bound.

For a cross with exponents $\alpha=s+i$, $\alpha'=s+k$ and
$A=\alpha+\alpha'-1$, the coordinate factor at order $n$ is
\[
 \frac{(A)_n}{A n!}\,
 \mathbb E[(T(1-T))^n]
 =\frac{(\alpha)_n(\alpha')_n}
 {n!(A+n)_{n+1}}=K_n(i,k),\quad
 T\sim\mathrm{Beta}(\alpha,\alpha').
\]
The logarithmic moment uses the tilted distribution
$\mathrm{Beta}(\alpha+n,\alpha'+n)$, giving the coordinate constant
\[
 \tfrac12\psi(n+1)-\tfrac12\psi(A+n)
 -\tfrac12[\psi(\alpha+n)+\psi(\alpha'+n)]
 +\psi(A+2n+1)=C_s/2+r_n(i,k).
\]
The digamma recurrence proves the rational expression for $r_n$.
The remainder of each $Q_{AB}$ after index $m$ is nonnegative and bounded
by $\log(1/b)(ABb^2)^{m+1}/(1-ABb^2)$. Monotonicity for $b\le a/4$
under \eqref{pf:cross-threshold} proves \eqref{pf:collected-overlap}.
We also have $B_n>0$, since $K_n$ has the Gram representation
\[
 K_n(i,k)=\frac{(s+i)_n(s+k)_n}{(n!)^2}
 \int_0^1 t^{\eta_0+n-1}(1-t)^n t^i t^k\,dt.
\]
It is positive definite on any finite collection of distinct monomials,
and its tensor square makes $B_n$ the squared norm of a nonzero
polynomial. Individual signed crosses can still be negative.

For the mass, let $T,U$ be independent $\mathrm{Beta}(\beta-1,1)$ variables.
Radial substitution gives
\begin{equation}\label{pf:polynomial-mass}
 A_D=\frac{\pi^2a^{-2}}{(\beta-1)^2}\mathcal A,
 \qquad \mathcal A=\mathbb E P(T,U)^p.
\end{equation}
For a rational positive Bernstein representation, first clear a common
denominator. Since scalar multiplication of $P$ cancels from $J_D$, normalize
it to $P=1+V$ with $|V|\le r<1$ by dividing by half the sum of the smallest
and largest Bernstein coefficient. Then, for $N\ge2$,
\begin{equation}\label{pf:polynomial-mass-tail}
 \left|\mathcal A-\sum_{k=0}^N\binom pk\mathbb EV^k\right|
 \le\frac{|\binom p{N+1}|r^{N+1}}{1-r}.
\end{equation}
Indeed the binomial series is uniformly convergent on $[-r,r]$ and the
absolute binomial coefficients decrease after index two, since
$1<p<2$. If $\beta-1=A/C$ is positive rational, the monomial moments
$\mathbb ET^j=A/(A+Cj)$ make every coefficient in the truncated sum
exactly rational. We calculate them by integer polynomial convolution
and contraction. Since the mass has a negative logarithmic coefficient
in the functional, its upper endpoint gives a lower bound for the score.

For polynomials with a large ratio of maximum to minimum Bernstein
coefficient, the following local version has a smaller tail. Partition
$[0,1]^2$ into rational rectangles, put $q=\beta-1>0$, and on
$\mathcal R=[l,r]\times[b,t]$ use coordinates
$x=(T-l)/(r-l)$ and $y=(U-b)/(t-b)$. Restrict the polynomial exactly and
express it in the Bernstein basis in $x,y$. If its coefficient bounds are
$0<m_{\mathcal R}\le M_{\mathcal R}$, set
\[
 h_{\mathcal R}=(m_{\mathcal R}+M_{\mathcal R})/2,\quad
 \rho_{\mathcal R}=
 \frac{M_{\mathcal R}-m_{\mathcal R}}{M_{\mathcal R}+m_{\mathcal R}}<1,
 \quad P=h_{\mathcal R}(1+W_{\mathcal R}).
\]
The Bernstein convex-hull property gives
$|W_{\mathcal R}|\le\rho_{\mathcal R}$ throughout the rectangle.
Define the unnormalized one-dimensional moments
\begin{equation}\label{pf:subrectangle-moments}
 \begin{split}
 M_i(l,r)&=\int_l^r qv^{q-1}
       \left(\frac{v-l}{r-l}\right)^i\,dv\\
 &=\frac q{(r-l)^i}\sum_{j=0}^i\binom ij(-l)^{i-j}
             \frac{r^{q+j}-l^{q+j}}{q+j}.
 \end{split}
\end{equation}
The identity follows by expanding $(v-l)^i$. The convention $0^q=0$
is valid because $q>0$. If
$W_{\mathcal R}^k=\sum c^{(k)}_{ij}x^iy^j$, put
\[
 T_{\mathcal R,N}=h_{\mathcal R}^{p}
 \sum_{k=0}^N\binom pk\sum_{i,j}c^{(k)}_{ij}M_i(l,r)M_j(b,t).
\]
Then for $N\ge2$,
\begin{equation}\label{pf:subrectangle-mass}
 \left|\mathcal A-\sum_{\mathcal R}T_{\mathcal R,N}\right|
 \le\sum_{\mathcal R}h_{\mathcal R}^{p}
 \frac{|\binom p{N+1}|\rho_{\mathcal R}^{N+1}}
      {1-\rho_{\mathcal R}}M_0(l,r)M_0(b,t).
\end{equation}
This follows by integrating the uniform binomial-tail bound against the
nonnegative density on each rectangle. Polynomial restriction,
conversion, and powers have rational coefficients and may be calculated
by integer convolution after clearing denominators. The endpoint powers
in \eqref{pf:subrectangle-moments} and $h_{\mathcal R}^p$ are enclosed
using outward-rounded logarithms and exponentials, since they need not
be rational. Exact rational algebra followed by interval contraction
preserves the bound when the moment formula has cancellation. For symmetric $P$
and the same partition on both axes, off-diagonal rectangles may be
counted twice and diagonal rectangles once. Otherwise all rectangles
must be included separately.

Combining \eqref{pf:polynomial-overlap} and
\eqref{pf:polynomial-mass} gives
\begin{equation}\label{pf:polynomial-functional}
 \begin{split}
 J_D\ge{}&\log b_0+2(1+\delta)\log(\beta-1)-2\delta\log\pi
 -(1+\delta)\log\mathcal A+2\delta\log a\\
 &+\log[\log(1/a)+\kappa_P-E_-/b_0].
 \end{split}
\end{equation}
More generally, if $Z_*>0$ is the lower expression in
\eqref{pf:collected-overlap} and $\mathcal A^*$ is an upper bound from
\eqref{pf:subrectangle-mass}, we obtain the bound used in the certificate,
\begin{equation}\label{pf:collected-functional}
 J_D\ge\log Z_*+2(1+\delta)\log(\beta-1)-2\delta\log\pi
       +2\delta\log a-(1+\delta)\log\mathcal A^*.
\end{equation}
Interval evaluation encloses this lower-bound expression. Its upper
endpoint need not bound $J_D$.
For rational $s$ and rational polynomial data all digamma shifts in
$\kappa_P$ reduce to rational corrections and the single constant
\[
 C_s=\frac2{2s-1}-2(s-1)^2
 \sum_{k\ge1}\frac1{k(k+s-1)(k+2s-2)}.
\]
The positive sum has tail at most $1/(2N^2)$ after $N$ terms. Hence
$C_s$ can be enclosed using a finite sum and this explicit remainder.
Together with the endpoint moments and
\eqref{pf:polynomial-fourier-contract}, these bounds give the
archimedean inputs to Proposition~\ref{fw:transfer}.
Section~\ref{sec:certificate} supplies the finite-place inputs and
evaluates the resulting margin.

We also use the following digamma enclosure. For $x>0$, integers $L\ge0$
and $m\ge1$, and $z=x+L$,
\[
 \left|\psi(x)-\left(\log z-\frac1{2z}
 -\sum_{k=1}^{m-1}\frac{B_{2k}}{2kz^{2k}}
 -\sum_{j=0}^{L-1}\frac1{x+j}\right)\right|
 \le\frac{|B_{2m}|}{2m z^{2m}},
\]
where $B_{2k}$ are rational Bernoulli numbers. To verify the remainder,
expand $(z^2+t^2)^{-1}$ as a finite geometric sum in Binet's identity
\cite[Equation~(5.9.15)]{DLMF},
\[
 \psi(z)=\log z-\frac1{2z}
 -2\int_0^\infty\frac{t}{(t^2+z^2)(e^{2\pi t}-1)}\,dt.
\]
The residual has a fixed sign and is bounded by the next integrated term,
$|B_{2m}|/(2m z^{2m})$. Apply the digamma recurrence $L$ times.
We evaluate the displayed expression with exact rational arguments and
Bernoulli numbers, outward arithmetic, and the explicit remainder.
The rational witness values are used without rounding the arguments.
